	\subsubsection{Construction of Teacher Variables}\label{sec:teachervar}
		This Section explains how we constructed the teacher variables that enter Table \ref{tab:teachereff} from the survey data that we collected among the Mathematics and Language teachers of the students in our sample.  \par
  
		
		\paragraph{Teacher effort.} For each teacher we observe the hours the teacher spends to prepare his/her classes, and the number of days the teacher was absent from school.\par
		
		\paragraph{Teacher's focus of instruction.} This variable measures whether the teacher is targeting his/her teaching to a specific part of the student ability distribution.\par
	
		For Mathematics and Language teachers separately we  construct a variable indicating the difficulty level at which the teacher is teaching using survey questions about how much of various components of the curriculum the teacher covered during the term, coupled with the teacher's assessment of the difficulty level of each component. For example, for Mathematics we present the teacher with a list of the 4 subfields taken from the official national curriculum (``Algebra and Functions", ``Geometry", ``Statistics and Probability", ``Trigonometry"), and for each subfield we present the teacher with a list of topics taken from the official national curriculum (for example, for ``Algebra and Functions" two topics are ``logarithmic and exponential function and analysis of their graphs" and ``solution of second degree equations"). In all, we presented Mathematics teachers with 13 topics and Language teachers with 11 topics. For each topic, we first ask the teacher what percentage he/she was able to cover during the first semester (which was over when the data collection started). Second, we ask the teacher to think of the average student in his/her $12^{th}$ grade class, and tell us whether he/she thinks that this student would find the topic easy or difficult to understand. The answers to these questions were collected as 5-point Likert scales. Finally, we multiply the coverage and difficulty within each Mathematics (Language) topic and sum over all topics. 
	